\scanpage{26}
\mspage{4}

\uline{Schémas abéliens polarisés principaux}.

Soit $A/S$ un schéma abélien, muni d'une section $\delta$ de $\underline{\mathrm{NS}}_{A/S}$, d'où un torseur $P^\delta=\underline{\mathrm{Pic}}^\delta_{A/S}$ sous le schéma abélien dual $B=\underline{\mathrm{Pic}}^0_{A/S}$, et un faisceau inversible $\mathscr L$ sur $A\times P^\delta$ \ldots. D'autre part, $\delta$ définit $\widetilde\delta\colon A\longrightarrow B$, et les opérations de $A$ sur $P^\delta$ via $\widetilde\delta$ sont celles déduites par transport de structure des opérations de $A$ sur lui-$=$ par translations.

Supposons que $\delta$ soit une polarisation, \deleted{\unsure{de|on}} alors on sait que les $H^i$ des $\mathscr L_x$ $(x\in P^\delta)$ sont nuls pour $i>0$, et que $\deleted{\unsure{R}}\operatorname{pr}_{2*}(\mathscr L)$ est donc un module loc libre de rang $\chi(\delta)$. Si $\chi(\delta)=1$, c'est un module inversible sur $P^\delta$, et il lui est associé un diviseur canonique $D_\delta$ \deleted{$\unsure{\Theta_\delta|\Theta}$} sur $A\times P^\delta$ relativement à $P^\delta$. D'autre part, $\widetilde\delta\colon A\mathbin{\unsure{-|\to}}B$ est alors un isom, et $P^\delta$ peut être regardé comme un $A$-torseur.

\deleted{\unsure{l'effet d'une translation}} \deleted{\unsure{$\Theta$}} Le diviseur $\deleted{\unsure{\Theta|D}}D_\delta$ est stable par opérations de $A$ opérant diagonalement sur $A\times P^\delta$ \deleted{\unsure{(sur les|Donc il}} Il provient donc d'un diviseur sur le quotient $(A\times P^\delta)/A=P^\delta$, noté \deleted{\unsure{aussi}} $\Theta_\delta$. Ce dernier a la vertu que $\chi(\Theta_\delta)=1$.

\uline{Proposition}. Soit $A$ un schéma abélien $/S$, \inserted{de dim relative $g$}, \inserted{Soit} $\delta$ une polarisation sur $A$ de degré 1. Alors il existe un $A$-torseur
% archive sheet number: 25
